\subsubsection{Random Forest}

Como hemos comentado sobre los Random Forest, estos construyen un conjunto o ''bosque'' de árboles de decisión donde cada árbol se entrena de forma ligeramente diferente, introduciendo aleatoriedad controlada.

\subsubsection*{Justificación del modelo}

Se seleccionó Random Forest como uno de los clasificadores base debido a su robustez frente al sobreajuste, su capacidad para manejar características heterogéneas y su interpretabilidad mediante el análisis de importancia de variables.

\subsubsection*{Configuración experimental}

El proceso experimental para Random Forest se estructuró en tres fases: (1) establecimiento de un modelo baseline, (2) validación cruzada para evaluar robustez, (3) optimización de hiperparámetros mediante búsqueda sistemática.

\paragraph{División de datos}
El conjunto de datos se particionó en entrenamiento (80\%) y prueba (20\%) utilizando muestreo estratificado para mantener la proporción de clases en ambos conjuntos. Esta estratificación es crucial cuando existe desbalance entre voces reales y generadas.

\paragraph{Modelo baseline}
Se estableció una configuración inicial con hiperparámetros por defecto para obtener una referencia de rendimiento:

\begin{itemize}
	\item \textbf{n\_estimators}: 100.
	\item \textbf{max\_depth}: 10.
	\item \textbf{min\_samples\_split}: 5.
	\item \textbf{min\_samples\_leaf}: 2.
	\item \textbf{max\_features}: sqrt.
	\item \textbf{random\_state}: 12.
	\item \textbf{class\_weight}: balanced.
\end{itemize}

\subsubsection*{Validación cruzada estratificada}

Para evaluar la robustez del modelo y evitar que los resultados dependan de una única partición de los datos, se realizó validación cruzada estratificada con 5 conjuntos (folds). Esto es dividir los datos en k conjuntos y entrenar un modelo con k-1 y evaluarlo con el restante, haciendo k entrenamientos (ya que cada conjunto se usa para evaluación exactamente una vez). Nosotros usamos el modelo baseline de Random Forest para aplicar la validación cruzada, mezclamos los datos de entrenamiento antes de instanciar el modelo de validación cruzada y, por último, evaluamos con accuracy, f1 y ROC-AUC.

La validación cruzada proporciona una estimación más realista del rendimiento esperado en datos no vistos, reportándose la media y la desviación estándar de las métricas.

\subsubsection*{Optimización de hiperparámetros}

Se realizó una búsqueda sistemática de hiperparámetros mediante \texttt{GridSearchCV} con validación cruzada de 5 folds, optimizando el F1-score. El espacio de búsqueda incluyó:

\begin{itemize}
	% ¿Explicación de para qué? Todos son para ver el valor óptimo
	\item \textbf{n\_estimators}: Probamos los valores 50, 100 y 200.
	\item \textbf{max\_depth}: Probamos los valores 5, 10, 15.
	\item \textbf{min\_samples\_split}: Probamos 2, 5 y 10.
	\item \textbf{min\_samples\_leaf}: Probamos 1, 2 y 4.
	\item \textbf{max\_features}: Probamos ``sqrt'' y ``log2''. Este último toma el logaritmo en base 2 del número de características.
\end{itemize}

Tras esto instanciamos el GridSearchCV (búsqueda de parámetros (search) probando todas las combinaciones (grid) con validación cruzada (CV)) poniendo un modelo base de Random Forest con el estado aleatorio 12 y el peso de clases balanceado, como antes. A dicho grid search le ponemos el grid de parámetros que comentamos arriba, la validación cruzada de 5-folds que hicimos antes, y la métrica a optimizar de f1. Por último entrenamos el modelo con los datos de entrenamiento sin escalar.